\begin{textsample}{POS Dim 2 – 2015\_c5 – Score -1.00 – 2015\_c5/enade\_2015\_tec\_des\_graf\_941cb2ecb1.txt}  \label{ex:f2_pos_282}
CAPÍTULO 5 DISTRIBUIÇÃO DOS CONCEITOS 5.1 PANORAMA NACIONAL DA DISTRIBUIÇÃO DOS CONCEITOS O Conceito Enade é um indicador de qualidade que avalia o desempenho dos estudantes a partir dos resultados do Enade.
Os valores possíveis vão de 1 ( pior situação ) a 5 ( melhor situação ).
As Unidades de Observação com apenas um ou sem nenhum concluinte participante não obtêm o Conceito Enade, ficando Sem Conceito ( SC ).
A Tabela 5.1 apresenta a quantidade e distribuição de cursos de Tecnologia em Design Gráfico participantes do Enade/2015, por faixa de conceito e Grande Região.
Mais uma vez, ressaltando que a diferença entre os cursos tabulados neste capítulo e no capítulo 2 são os cursos sem conceito, em princípio, aqueles sem alunos concluintes que participassem da prova ou que tivessem tão somente um aluno concluinte¹⁸.
Observando -se os dados da Tabela 5.1, nota -se que, dos 77 cursos participantes, 38 ( 49,4 \% ) classificaram -se com conceito 3, o valor modal.
Esse foi também o conceito modal em mais três regiões : Nordeste ( 66,7 \% ), Sudeste ( 55,8 \% ) e Sul ( 33,3 \% ).
Na região Centro-Oeste, a moda foi o conceito 4 ( 33,3 \% ), enquanto na Norte, a moda foi o conceito 2 ( 50,0 \% ).
O conceito 2 foi o segundo mais frequente em nível nacional ( 23,4 \%, correspondendo a 18 cursos ), e o conceito 4, o terceiro ( 15,6 \%, correspondendo a 12 cursos ).
Houve, ainda, quatro cursos ( 5,2 \% ) que receberam conceito 5, e dois cursos ( 2,6 \% ) que receberam conceito 1.
Dos 77 cursos de Tecnologia em Design Gráfico, três cursos ficaram sem conceito ( SC ). ¹⁸ Estes cursos com somente um concluinte não foram considerados no capítulo 2.
A região Norte participou com quatro cursos ou 5,2 \% do total nacional.
Nesta região, o conceito 2 foi atribuído a dois cursos dos quatro participantes, o que equivale a 50,0 \% do total regional, além de ser o conceito modal na região.
Um curso ( 25,0 \% ) obteve conceito 3.
Além disso, um curso ficou sem conceito ( SC ) enquanto nenhum outro recebeu os conceitos 1, 4 ou 5.
A região Nordeste participou com 12 cursos ou 15,6 \% do total nacional.
Nessa região, 8 cursos, 66,7 \% em termos regionais, obtiveram conceito 3, o conceito modal para a região.
Os conceitos 2 e 4 foram atribuídos a dois cursos ( 16,7 \% ), cada um.
Nenhum curso apresentou os conceitos 1 ou 5, bem como nenhum curso ficou sem conceito ( SC ).
Dos 43 cursos participantes da região Sudeste ( 55,8 \% do total nacional ), o conceito 3 também foi o conceito modal, e foi obtido por 24 ( 55,8 \% ) cursos, seguido do conceito 2, atribuído a dez cursos, ou 23,3 \%.
Sete cursos ( 16,3 \% ) receberam o conceito 4.
Não houve atribuição dos conceitos 1 ou 5, enquanto dois cursos ficaram sem conceito ( SC ).
Os 12 cursos da região Sul ( 15,6 \% do total nacional ) ficaram distribuídos em todas as faixas de conceitos.
Há predominância do conceito 3 com 33,3 \% ( conceito modal ), ou seja, quatro dos 12 cursos participantes na região Sul.
O conceito 2 foi atribuído a três cursos ( 25,0 \% ), assim como o foi o conceito 5, com os mesmos 25,0 \%.
Os conceitos 1 e 4 foram atribuídos a 1 curso ( 8,3 \% ) cada um.
Conforme já se disse, nenhum curso ficou sem conceito ( SC ).
Já dos seis cursos participantes na região Centro-Oeste ( 7,8 \% do total nacional ), dois ( 33,3 \% ) receberam conceito 4, o conceito modal.
Os demais cursos receberam, cada qual, um conceito diverso : conceitos 1, 2, 3 e 5, cada um deles correspondendo a 16,7 \% do total regional.
Nenhum curso ficou sem conceito ( SC ) nessa região.
O Gráfico 5.1 apresenta a distribuição do Conceito Enade segundo Grande Região : poligonais mais à esquerda representam Grandes Regiões com pior distribuição, e poligonais mais à direita, Grandes Regiões com melhores conceitos.
A região Norte ( linha azul ) apresenta os cursos em IES com conceitos mais baixos, ao passo que as regiões Centro-Oeste ( linha preta ) e Sul ( linha roxa ) apresentam as distribuições com valores maiores.
Cursos em IES das regiões Sudeste ( linha vermelha ) e Nordeste ( linha verde ) apresentam uma situação intermediária. 5.2 CONCEITOS POR CATEGORIA ADMINISTRATIVA E POR GRANDE REGIÃO A Tabela 5.2 apresenta a distribuição dos cursos participantes do Enade/2015 de Tecnologia em Design Gráfico, por Categoria Administrativa, de acordo com os conceitos por eles alcançados, segundo as Grandes Regiões brasileiras.
Dos 77 cursos participantes, 5 ( 6,5 \% ) eram ministrados em Instituições Públicas e 72 ( 93,5 \% ), em Instituições Privadas.
De acordo com as informações da Tabela 5.2, em termos nacionais, entre Instituições Públicas e Privadas, das quatro IES que receberam o conceito 5, duas eram Públicas e duas Privadas.
Dos cinco cursos participantes de IES Públicas, os conceitos 4 e 5 foram os valores modais, atribuídos a dois cursos ( 40,0 \% ) cada um.
Um curso ( 20,0 \% ) recebeu conceito 1.
Entre as Instituições Públicas participantes, a nenhum curso foram atribuídos os conceitos 2 ou 3, tampouco houve curso sem conceito ( SC ).
Na rede privada, o conceito modal foi 3, com 38 cursos ou 52,8 \% da categoria.
Entre os demais cursos participantes, apenas um ( 1,4 \% ) recebeu conceito 1, e 18 ( 25,0 \% ) receberam conceito 2.
O conceito 4 foi atribuído a dez cursos ( 13,9 \% ), e o conceito 5, a dois cursos ( 2,8 \% ).
Nesta Categoria Administrativa, outros três cursos ficaram sem conceito ( SC ).
Na análise por região, observa -se que, na região Norte, as Instituições Privadas participaram com os quatros cursos ( 100,0 \% da região ), dos quais um ficou sem conceito ( SC ) enquanto nenhum ficou com conceitos 1, 4 ou 5.
O conceito modal para as Instituições Privadas na região foi 2, com dois cursos, correspondendo a 50,0 \%.
E o terceiro curso foi alocado no conceito 3, perfazendo os 25,0 \% restantes.
As Instituições Públicas não participaram com curso na região Norte.
Na região Nordeste, a rede privada concentrou quase a totalidade dos 12 cursos participantes, o equivalente a 91,7 \% do total da região : foram 11 cursos.
Dentre os cursos oferecidos por IES Privadas no Nordeste, oito cursos ficaram alocados no conceito 3, que foi a moda da região, correspondente a 72,7 \%.
Os demais conceitos foram assim distribuídos : dois cursos obtiveram conceito 2 ( 18,2 \% ) ; um único curso obteve conceito 4 ( 9,1 \% ).
Nenhum curso ficou sem conceito ( SC ) nesta combinação de categoria e região, tampouco foram atribuídos os conceitos 1 ou 5.
A única Instituição Pública dessa região obteve o conceito 4.
Na região Sudeste, a proporção de cursos da rede privada, 97,7 \%, foi também bastante acentuada, abaixo apenas da região Norte, que, conforme já se disse, tem todos os seus cursos nessa categoria.
Dos 43 cursos oferecidos, 42 são da rede privada.
Nessa categoria, na região Sudeste, o conceito modal foi 3 ( 24 cursos, correspondendo a 57,1 \% ), seguido do conceito 2, atribuído a 23,8 \% dos cursos, isto é, dez instituições.
O conceito 4 foi atribuído a seis cursos ( 14,3 \% ).
Dois cursos ficaram sem conceito ( SC ) nessa combinação de categoria e região, enquanto a nenhum foi atribuído os conceitos 1 ou 5.
A única Instituição Pública dessa região obteve o conceito 4.
O predomínio das Instituições Privadas também se verifica nos cursos participantes da região Sul, 83,3 \% do total regional, ou seja, dez dos 12 cursos ali oferecidos.
Desse total, quatro cursos obtiveram o conceito modal 3 ( 40,0 \% ) ; três receberam conceito 2 ( 30,0 \% ).
Os conceitos 1, 4 e 5 foram atribuídos a um curso ( 10,0 \% ) cada um.
Nenhum curso ficou sem conceito ( SC ) nesta combinação de categoria e região.
As duas Instituições Públicas participantes obtiveram conceito 5.
Na região Centro-Oeste, cinco dos seis cursos participantes eram de Instituições Privadas ( 83,3 \% em termos regionais ).
Destes, dois concentraram -se no conceito 4, conceito modal.
Os conceitos 2, 3 e 5 foram atribuídos a um curso ( 20,0 \% ) cada um.
Nenhum curso obteve conceito 1 ou ficou sem conceito ( SC ) nessa combinação de categoria e região.
A única Instituição Pública dessa região obteve o conceito 1.
O Gráfico 5.2 apresenta a distribuição do Conceito Enade segundo a Categoria Administrativa da IES.
Os cursos em IES Públicas ( linha azul ) apresentam uma poligonal mais à direita do que os em IES Privadas ( linha verde ) e, consequentemente, uma distribuição de Conceitos Enade com valores maiores. 5.3 CONCEITOS POR ORGANIZAÇÃO ACADÊMICA E POR GRANDE REGIÃO Na Tabela 5.3, encontra -se a distribuição dos conceitos atribuídos aos cursos participantes do Enade/2015 na Área de Tecnologia em Design Gráfico, por Organização Acadêmica, segundo as Grandes Regiões brasileiras.
Dos 77 cursos de Tecnologia em Design Gráfico participantes, 40 eram oferecidos em Universidades ; 19 cursos em Centros Universitários ; 18, em Faculdades.
Essa distribuição corresponde a, respectivamente, 51,9 \%, 24,7 \% e 23,4 \% dos cursos.
De acordo com os dados apresentados, dos quatro cursos avaliados com conceito 5, três eram vinculados a Universidades.
Dois cursos oferecidos em Universidades ficaram sem conceito ( SC ).
Esse tipo de Organização Acadêmica teve o conceito 3 como modal, com 16 cursos ( 40,0 \% ).
Os demais cursos de Universidades avaliados receberam os conceitos 1 ( dois cursos ), 2 ( 12 cursos ), 4 ( cinco cursos ) e conceito 5 ( três cursos, como já mencionado ).
Entre os cursos em Centros Universitários, o conceito modal também foi 3, com 12 cursos ( 63,2 \% ).
Nesse tipo de Organização Acadêmica, nenhum curso ficou sem conceito ( SC ) ou obteve conceito 1, e os outros cursos receberam os conceitos 2 ( três cursos ), 4 ( três cursos ), 5 ( apenas um curso ).
Nas Faculdades, apenas um curso ficou sem conceito ( SC ), e dez ( 55,6 \% ) ficaram alocados no conceito 3, conceito modal, seguido de quatro cursos que obtiveram conceito 4 e de três cursos que receberam conceito 2.
Nesse tipo de Organização Acadêmica, nenhum curso obteve conceitos 1 e 5.
Considerando -se separadamente as regiões brasileiras, na região Norte, foram oferecidos apenas quatro cursos.
Conforme se pode constatar na Tabela 5.3, as Universidades não oferecem curso.
Os Centros Universitários foram representados por Na região Nordeste, as Universidades participaram com apenas dois dos 12 cursos na Área de Tecnologia em Design Gráfico, aos quais foram atribuídos os conceitos 2 e 4.
Os Centros Universitários contaram com somente um curso participante na região Nordeste, e este obteve conceito 4.
As Faculdades foram representadas por nove cursos na região Nordeste, dos quais oito ( 88,9 \% ) receberam o conceito modal 3, e o outro ficou com o conceito 2.
Na região Sudeste, as Universidades concentraram 28 dos 43 cursos da região.
Entre os cursos em Universidades, nessa região, o conceito modal foi 3 com 14 cursos, seguido do conceito 2, com nove cursos, e do conceito 4, com três.
Nesse tipo de Organização Acadêmica, dois cursos ficaram sem conceito ( SC ) enquanto nenhum obteve conceitos 1 ou 5.
Os Centros Universitários participaram com nove cursos na região Sudeste, dos quais oito obtiveram o conceito modal 3, e um curso obteve conceito 4.
Nenhum curso recebeu conceitos 1, 2 ou 5 nessa categoria e região, tampouco ficou sem conceito ( SC ).
As Faculdades foram representadas por seis cursos na região Sudeste dos quais a metade recebeu o conceito modal 4.
Dois cursos obtiveram conceito 3 ; um curso obteve conceito 2.
Portanto, nesse tipo de Organização Acadêmica, nenhum curso ficou sem conceito ( SC ) nem obteve conceitos 1 ou 5.
Dos 12 cursos da região Sul, oito eram de Universidades, para os quais o conceito modal foi 5, com três cursos.
Nesse tipo de organização, nenhum dos cursos ficou sem conceito ( SC ).
Os demais receberam o conceito 3 ( dois cursos ) e os conceitos 1, 2 e 4 ( um curso cada um ).
Os quatro cursos participantes de Centros Universitários da região Sul obtiveram apenas os conceitos 2 e 3, distribuídos por dois cursos cada um.
Não houve cursos da área de Tecnologia em Design Gráfico oferecidos pelas Faculdades, na região Sul.
Na região Centro-Oeste, dos seis cursos oferecidos, apenas dois eram de Universidades, que receberam os conceitos 1 e 2 cada um.
Os Centros Universitários da região Centro-Oeste contaram com três cursos, que obtiveram, cada um deles, um conceito diferente : 3, 4, e 5.
Nenhum curso ficou sem conceito ( SC ) ou recebeu os conceitos restantes na categoria.
O único curso oferecido por uma Faculdade na região Centro-Oeste obteve o conceito 4.
O Gráfico 5.3 apresenta a distribuição do Conceito Enade segundo a Organização Acadêmica da IES.
Os cursos em Universidades ( linha azul ) apresentam uma poligonal mais à esquerda do que as demais, denotando uma pior distribuição de Conceitos Enade.
Mais à direita, temos os cursos oferecidos em Faculdades ( linha vermelha ) e as IES de Centros Universitários ( linha verde ), cujas poligonais em muitos pontos praticamente se superpõe.
Por estarem mais à direita, apresentam, consequentemente, uma distribuição de Conceitos Enade com maiores valores, embora estas últimas tenham suas poligonais também se entrecruzando com a de Universidades.

% matched lemmas: 
\end{textsample}
